%=====================================================================
\chapter{Answers to Checkpoints}\label{app:answers}
%=====================================================================
\normalfigures

Answers to the three-question Checkpoint box in each chapter. Attempt them before
reading. Where an answer requires judgement, the reasoning matters more than the
conclusion.

\newcommand{\ansch}[1]{\subsection*{\color{primary}Chapter #1}}
\newenvironment{ans}{\begin{enumerate}[leftmargin=*, itemsep=4pt, topsep=3pt,
  label=\textbf{\arabic*.}]}{\end{enumerate}}

%---------------------------------------------------------------------
\ansch{1 --- What Machine Learning Is}
\begin{ans}
  \item No. The rule was written by the owner and the thermostat's performance does
        not improve with experience: there is no $E$. It is automation, not learning.
  \item (a) regression; (b) clustering (unsupervised); (c) classification;
        (d) reinforcement learning.
  \item Because for any unseen instance, every hypothesis consistent with the data
        that calls it positive has a twin, equally consistent, that calls it negative.
        The data says nothing about unseen instances; only an assumption --- a bias
        --- can.
\end{ans}

\ansch{2 --- Mathematics for Machine Learning}
\begin{ans}
  \item $\sqrt{(4-1)^2 + (6-2)^2} = \sqrt{9 + 16} = 5$.
  \item $L'(w) = 2(w + 1) = 6$ at $w = 2$, so $w \leftarrow 2 - 0.25 \times 6 = 0.5$. The
        minimum is at $-1$: the gap fell from 3 to 1.5, so yes, closer.
  \item $H = -(0.5\log_2 0.5 + 2 \times 0.25\log_2 0.25) = 0.5 + 0.5 + 0.5 = 1.5$ bits.
\end{ans}

\ansch{3 --- The Machine Learning Workflow and Toolkit}
\begin{ans}
  \item At prediction time only the training data is known. Statistics computed with
        test rows leak information about the test set into training and make the
        evaluation optimistic; new data must be transformed with the stored training
        statistics.
  \item One-hot encode it into three 0/1 columns. Codes 1, 2, 3 would tell a model that
        ICMP is ``three times'' TCP and that UDP lies between them --- an order and
        distances that do not exist.
  \item The majority-class baseline: if 97\% of machines never fail, predicting ``no
        failure'' is 97\% accurate. Then ask for recall on the failures.
\end{ans}

\ansch{4 --- Concept Learning and the Version Space}
\begin{ans}
  \item Yes: any day satisfying $\langle$Sunny, Warm, ?, Strong, ?, ?$\rangle$ also satisfies
        $\langle$Sunny, ?, ?, Strong, ?, ?$\rangle$. The reverse is false: a sunny, cold,
        windy day satisfies the first but not the second.
  \item $\langle$Sunny, Warm, ?, Strong, Warm, Same$\rangle$ --- the hypothesis after Day 2,
        since Day 3 is negative and ignored.
  \item For each concept in its version space that labels the unseen instance positive,
        there is another, identical except for that instance, labelling it negative;
        both are consistent with the training data. The votes are always equal.
\end{ans}

\ansch{5 --- Linear Regression}
\begin{ans}
  \item $\bar{x} = 1$, $\bar{y} = 3$; $w = \frac{(-1)(-2) + 0 + (1)(2)}{1 + 0 + 1} = 2$;
        $b = 3 - 2 = 1$. So $\hat{y} = 2x + 1$, which passes through all three points.
  \item The model's squared error on the test set is 20\% larger than that of simply
        predicting the training mean: it is worse than the baseline.
  \item Because $\hat{y}$ is a weighted sum of its features --- here $x$ and $x^2$ --- so
        it is linear in the weights, and least squares applies unchanged.
\end{ans}

\ansch{6 --- Logistic Regression}
\begin{ans}
  \item $p = 1/(1 + e^{-1.2}) = 1/1.301 = 0.769$; odds $= e^{1.2} = 3.32$ (about 10 to 3).
  \item $e^{0.7} = 2.01$: the odds roughly double.
  \item A positive predicted at 0.4 costs $-\ln 0.4 = 0.916$; a negative predicted at
        0.4 costs $-\ln 0.6 = 0.511$. The positive costs more, because 0.4 is further
        from its true label.
\end{ans}

\ansch{7 --- Overfitting, Bias--Variance and Regularisation}
\begin{ans}
  \item Variance (overfitting). Remedies: more data, stronger regularisation or a
        simpler model, fewer features, early stopping, an averaging ensemble.
  \item Bias$^2 = (2 - 3)^2 = 1$. Variance $= 0$: every prediction is identical.
  \item The L1 constraint region is a diamond with corners on the axes; the error
        contours usually first touch it at a corner, where some weights are exactly
        zero. The L2 region is a circle with no corners, so weights shrink but rarely
        reach zero.
\end{ans}

\ansch{8 --- Measuring Classifier Performance}
\begin{ans}
  \item TP $= 100$, FP $= 300$, FN $= 20$. Precision $= 100/400 = 0.25$; recall
        $= 100/120 = 0.83$.
  \item The precision--recall curve, or precision at the recall the application needs,
        and the number of alerts per day. With 0.1\% positives a high ROC-AUC is
        compatible with most alerts being false.
  \item $e = 0.15$: $1.96\sqrt{0.15 \times 0.85/100} = 0.070$, so roughly 8\% to 22\%.
\end{ans}

\ansch{9 --- Decision Tree Learning}
\begin{ans}
  \item The parent has entropy 1. Each child has entropy
        $-0.75\log_2 0.75 - 0.25\log_2 0.25 = 0.811$ and holds half the examples, so
        Gain $= 1 - 0.811 = 0.189$.
  \item Every ID is unique, so splitting on it gives one pure leaf per student: maximum
        gain, no generalisation. Use the gain ratio, whose split-information term
        penalises many-valued attributes --- or remove identifiers.
  \item A restriction bias limits which hypotheses can be represented
        (\textsc{Candidate-Elimination}'s conjunctions); a preference bias searches a
        complete space but prefers some hypotheses. ID3 has a preference bias: shorter
        trees, high-gain attributes near the root.
\end{ans}

\ansch{10 --- Bayesian Learning and Naive Bayes}
\begin{ans}
  \item Disease: $0.9 \times 0.01 = 0.009$; healthy: $0.05 \times 0.99 = 0.0495$. So
        $h_{MAP}$ is \emph{healthy}; $h_{ML}$ is \emph{disease} ($0.9 > 0.05$). The
        posterior probability of disease is $0.009/0.0585 = 0.15$.
  \item That the attributes are independent of each other given the class. When they
        are correlated the same evidence is counted several times, and the
        probabilities are pushed towards 0 and 1 --- over-confident --- even though the
        ranking of classes is often still right.
  \item The product of hundreds of small probabilities underflows to zero in
        floating-point arithmetic. Logarithms turn the product into a sum of
        manageable numbers.
\end{ans}

\ansch{11 --- Instance-Based Learning: k-Nearest Neighbours}
\begin{ans}
  \item Distances from (50, 55): F 11.18 (Fail), D 15.81 (Fail), C 22.36 (Pass), E
        23.77, B 33.54, A 50.0. The three nearest vote Fail, Fail, Pass: \textbf{Fail}.
  \item It does no work until asked to classify, and then searches the stored examples.
        The cost is paid at prediction time --- a distance to every stored example for
        each query --- and all the training data must be kept.
  \item Bias rises and variance falls. At $k = 1$ the boundary follows every example; at
        $k = n$ every query gets the majority class.
\end{ans}

\ansch{12 --- Support Vector Machines}
\begin{ans}
  \item $\|\mathbf{w}\| = 5$, so the margin is $2/5 = 0.4$. $f(1, 1) = 3 + 4 - 5 = 2 > 0$:
        positive.
  \item The margin narrows and the model tolerates fewer violations, so training
        errors fall --- and the risk of overfitting rises.
  \item Because both the SVM's optimisation and its decision function use the training
        examples only through dot products $\mathbf{x}_i \cdot \mathbf{x}_j$, and each dot
        product can be replaced by a kernel.
\end{ans}

\ansch{13 --- Artificial Neural Networks}
\begin{ans}
  \item $(0, 1)$: $-1 + 0 + 1 = 0$, not greater than 0, so output 0. $(1, 0)$:
        $-1 + 2 + 0 = 1 > 0$, so output 1.
  \item XOR's positive and negative points lie on opposite diagonals; no single line
        separates them. A hidden layer with one unit for OR and one for NAND, combined
        by an AND output unit, represents it.
  \item A hidden unit affects the error only through the output units it feeds, so
        by the chain rule its share of the blame is the output deltas weighted by its
        connections to them: $\delta_h = o_h(1 - o_h)\sum_k w_{kh}\delta_k$.
\end{ans}

\ansch{14 --- Ensemble Learning}
\begin{ans}
  \item $3(0.8^2)(0.2) + 0.8^3 = 0.384 + 0.512 = 0.896$.
  \item $\alpha = \tfrac{1}{2}\ln(0.8/0.2) = \tfrac{1}{2}\ln 4 = 0.693$. Misclassified
        examples are multiplied by $e^{0.693} = 2$, correct ones by $e^{-0.693} = 0.5$,
        before renormalising.
  \item A tree is unstable: different bootstrap samples give very different trees, whose
        errors partly cancel when averaged. Logistic regression is stable: models on
        different bootstrap samples are nearly identical, so averaging them changes
        little.
\end{ans}

\ansch{15 --- Clustering}
\begin{ans}
  \item 1 and 4 go to 2; 6, 7 and 9 go to 8. New centroids: $(1 + 4)/2 = 2.5$ and
        $(6 + 7 + 9)/3 = 7.33$.
  \item Neither step can increase the WCSS and there are finitely many possible
        assignments, so it must stop. But it stops at a local minimum that depends on
        the starting centroids.
  \item $s = (2 - 3)/3 = -0.33$. The point is, on average, closer to another cluster
        than to its own: it is probably in the wrong cluster.
\end{ans}

\ansch{16 --- Dimensionality Reduction and Self-Organizing Maps}
\begin{ans}
  \item $(6 + 3)/(6 + 3 + 1) = 0.9$.
  \item PCA seeks directions of large variance, and a feature measured in large units
        has large variance whether or not it is informative; it would dominate the
        components purely because of its units.
  \item The BMU is the single unit whose weights are closest to the input
        (competition). Its neighbourhood is the units near it \emph{on the grid}.
        Updating the neighbours as well makes adjacent units similar, which produces a
        smooth, topology-preserving map.
\end{ans}

\ansch{17 --- Semi-Supervised Learning and the EM Algorithm}
\begin{ans}
  \item Examples in the same cluster tend to share a label, so the decision boundary
        should pass through low-density regions. It fails when classes overlap inside
        one cluster --- for example fraudulent transactions designed to look exactly
        like normal ones.
  \item The E-step computes, with the current parameters, each example's
        responsibilities (the probability that each component produced it). The M-step
        re-estimates the parameters from the responsibility-weighted examples.
  \item $k$-means makes hard assignments and estimates only means (implicitly round,
        equal clusters); EM for a Gaussian mixture makes soft assignments and also
        estimates covariances and mixing weights, maximising the likelihood.
\end{ans}

\ansch{18 --- Hidden Markov Models}
\begin{ans}
  \item D$\to$H$\to$H $= 0.3 \times 0.8 = 0.24$; D$\to$D$\to$H $= 0.7 \times 0.3 = 0.21$;
        total $0.45$.
  \item $\alpha_t(i)$ is the probability of the observations up to $t$ \emph{and} being
        in state $i$, summed over every path; $\delta_t(i)$ is the probability of the
        single best path ending in $i$ --- a maximum instead of a sum.
  \item There are $N^T$ state sequences, exponentially many. The forward algorithm
        stores $\alpha_t$ for each state and reuses it at the next step, which costs
        $N^2 T$.
\end{ans}

\ansch{19 --- Markov Decision Processes}
\begin{ans}
  \item $4 + 0.5 \times 4 + 0.25 \times 4 = 7$.
  \item $V^{\pi}(s)$ is the expected return from $s$ following $\pi$; $Q^{\pi}(s, a)$ is the
        expected return from taking $a$ first and then following $\pi$. So
        $V^{\pi}(s) = Q^{\pi}(s, \pi(s))$.
  \item Because a machine that is operated while degraded eventually fails, and a failed
        machine that keeps being ``operated'' earns nothing for ever. The small immediate
        gain of operating is outweighed by the lost future earnings that repair
        preserves.
\end{ans}

\ansch{20 --- Reinforcement Learning}
\begin{ans}
  \item $(10 + 4 + 7)/3 = 7$; after a return of 3, $7 + (3 - 7)/4 = 6$.
  \item An arm that happened to pay badly the first time is never tried again, so its
        low estimate is never corrected; greedy selection can lock onto a worse arm for
        ever.
  \item $Q \leftarrow 4 + 0.5\,(1 + 0.9 \times 6 - 4) = 4 + 0.5 \times 2.4 = 5.2$.
\end{ans}

\ansch{21 --- Responsible and Ethical Machine Learning}
\begin{ans}
  \item $30/45 = 0.67 < 0.8$: it fails the four-fifths rule.
  \item Other features act as proxies and let the model reconstruct the attribute, and
        the labels themselves may encode past unfairness.
  \item Equal opportunity; it compares true positive rates (recall) across groups.
\end{ans}

\ansch{22 --- The Machine Learning Project}
\begin{ans}
  \item For example: ``On a freshly labelled sample of 5{,}000 messages from the test
        week, the filter blocks at least 95\% of spam while blocking fewer than 1 in
        1{,}000 legitimate messages.''
  \item They differ by 0.007, far less than one standard deviation of either: they are
        indistinguishable. Choose on other grounds --- simplicity, speed,
        explainability, fairness.
  \item Any three of: the random seeds; the data version and the split indices; the
        code version (commit); the library versions; the hyperparameters.
\end{ans}
